\sgasection{5}{First corollaries of the representability theorem}
Let $H,G$ be as in 4.2. Put
\[
M=\uHom_{S\text{-gr}}(H,G),
\]
which is a smooth $S$-prescheme, separated over $S$, by 4.2. Note that $G$ acts on the functor $M=\uHom_{S\text{-gr}}(H,G)$ by
\[
(g,u)\mto\operatorname{int}(g)\circ u,
\]
whence a canonical morphism
\begin{equation}\tag{x}
G\times_S M\to M\times_S M,
\end{equation}
whose components are the preceding morphism $G\times_S M\to M$, and the second projection $G\times_S M\to M$.
\begin{corollary}{5.1}\label{XI.5.1}
The preceding morphism (x) is smooth.
\end{corollary}
This follows from 4.2 and 2.3. This statement is equivalent to the following:
\begin{corollary}{5.2}\label{XI.5.2}
Let $u_1,u_2:H\rightrightarrows G$ be two homomorphisms of $S$-groups. Then the subfunctor $\uTransp(u_1,u_2)$ of $G$ (\cf 2.4) is representable by a closed sub-prescheme of $G$, smooth over $S$.
\end{corollary}
It remains only to prove that $\uTransp(u_1,u_2)\to G$ is indeed a closed immersion, which follows from the fact that it is the kernel of a pair of morphisms $G\rightrightarrows M$ (namely $g\mto\operatorname{int}(g)u_1$ and the ``constant'' morphism $g\mto u_2$), and from the fact that $M$ is separated over $S$. In particular:
\begin{corollary}{5.3}\label{XI.5.3}
Let $u:H\to G$ be a morphism of $S$-groups. Then $\uCentr_G(u)=\uTransp(u,u)$ is representable by a closed group sub-prescheme of $G$, smooth over $S$. In addition,
\[G/\uCentr_G(u)\]
is representable by an open sub-prescheme of $M$.
\end{corollary}
It remains to prove this last point. Now the morphism
\[
g\mto\operatorname{int}(g)\circ u
\]
from $G$ to $M$ is smooth of finite type by 5.2, hence is an open morphism (EGA~IV~6.6), and if $U$ denotes its image, with the structure induced by $M$, the induced morphism $G\to U$ is smooth, surjective, of finite type, hence covering for the faithfully flat quasi-compact topology. Moreover, it is obvious that the preceding morphism $G\to M$ makes $G$ a formally principal homogeneous sheaf under $\uCentr_G(u)_M$, which implies that the sheaf $G/\uCentr_G(u)$ is indeed representable by $U$.

\begin{corollary}{5.4}\label{XI.5.4}
Let $u_1,u_2:H\rightrightarrows G$ be two homomorphisms of $S$-groups. The following conditions are equivalent:
\begin{enumeratei}
\item $u_1$ and $u_2$ are locally conjugate for the étale topology.
\item[(i bis)] $u_1,u_2$ are locally conjugate in the sense of the faithfully flat quasi-compact topology.
\item For every $s\in S$, denoting by $\overline s$ the spectrum of an algebraic closure of $\kappa(s)$, the morphisms $u_{1,\overline s},u_{2,\overline s}:H_{\overline s}\rightrightarrows G_{\overline s}$ are conjugate by an element of $G(\overline{\kappa(s)})$.
\item[(ii bis)] The structural morphism $\uTransp(u_1,u_2)\to S$ is surjective.
\item $\uTransp(u_1,u_2)$ is a torsor under the action of the smooth $S$-group prescheme of finite type $\uCentr_G(u_1)$.
\end{enumeratei}
\end{corollary}
(i) $\Rightarrow$ (i bis) and (ii) $\Rightarrow$ (ii bis) are trivial (the second thanks to the Nullstellensatz); moreover (i bis) $\Rightarrow$ (ii) by the ``finite extension principle'' (EGA~IV~9). On the other hand, (ii bis) $\Rightarrow$ (iii) thanks to the fact that $\uTransp(u_1,u_2)$ is smooth over $S$, hence flat over $S$, and of finite type, hence quasi-compact over $S$, so it is faithfully flat quasi-compact over $S$ if and only if its structural morphism is surjective. Since, on the other hand, it is formally principal homogeneous under $\uCentr_G(u_1)$, which is faithfully flat quasi-compact over $S$, one sees that this last condition is also equivalent to saying that $\uTransp(u_1,u_2)$ is a torsor under $\uCentr_G(u_1)$ (understood: in the sense of the faithfully flat quasi-compact topology). Finally (iii) $\Rightarrow$ (i) thanks to the fact that $\uTransp(u_1,u_2)$ is smooth over $S$ and to ``Hensel's lemma'' in the form 1.10.
\begin{remark}{5.5}\label{XI.5.5}
For fixed $u_1=u:H\to G$, the functor $\Sch^\circ_{/S}\to\Ens$ that associates to every $T$ over $S$ the set of homomorphisms of $T$-groups $u_2:H_T\to G_T$ that are conjugate to $u_T:H_T\to G_T$ locally for the étale topology is precisely representable by the open subset of $M$, isomorphic to $G/\uCentr_G(u)$, considered in 5.3.
\end{remark}

Let us sketch the variants of the preceding results obtained by applying 4.1 instead of 4.2. Thus let $G$ be a group prescheme smooth and affine over $S$, and now denote by $M$ the smooth $S$-prescheme, separated over $S$, representing the functor considered in 4.1. One again has actions of $G$ on $M$:
\[
(g,H)\mto\operatorname{int}(g)(H),
\]
whence, as above, a morphism
\begin{equation}\tag{x bis}
G\times_S M\to M\times_S M.
\end{equation}
\begin{corollary}{5.1 bis}\label{XI.5.1bis}
The preceding morphism is smooth.
\end{corollary}
Follows from 4.1 and 2.3 bis. One again concludes:
\begin{corollary}{5.2 bis}\label{XI.5.2bis}
Let $H_1,H_2$ be two subgroups of multiplicative type of $G$. Then the subfunctor $\uTransp_G(H_1,H_2)$ of $G$ is representable by a closed sub-prescheme of $G$, smooth over $S$.
\end{corollary}
In particular:
\begin{corollary}{5.3 bis}\label{XI.5.3bis}
Let $H$ be a subgroup of multiplicative type of $G$. Then the group subfunctor $\uNorm_G(H)$ of $G$ is representable by a closed sub-prescheme of $G$, smooth over $S$. In addition, the quotient $G/\uNorm_G(H)$ is representable by an open sub-prescheme of $M$.
\end{corollary}
\begin{corollary}{5.4 bis}\label{XI.5.4bis}
Let $H_1,H_2$ be two subgroups of multiplicative type of $G$. The following conditions are equivalent:
\begin{enumeratei}
\item $H_1$ and $H_2$ are locally conjugate for the étale topology.
\item[(i bis)] $H_1$ and $H_2$ are locally conjugate for the faithfully flat quasi-compact topology.
\item For every $s\in S$, denoting by $\overline s$ the spectrum of an algebraic closure of $\kappa(s)$, the subgroups $H_{1,\overline s},H_{2,\overline s}$ of $G_{\overline s}$ are conjugate by an element of $G(\overline{\kappa(s)})$.
\item[(ii bis)] The structural morphism $\uTransp(H_1,H_2)\to S$ is surjective.
\item $\uTransp(H_1,H_2)$ is a principal homogeneous bundle under the action of the smooth $S$-group prescheme of finite type $\uNorm_G(H_1)$.
\end{enumeratei}
\end{corollary}
\begin{remark}{5.5 bis}\label{XI.5.5bis}
Remark 5.5 also carries over to the present case.
\end{remark}
\begin{remark}{5.6}\label{XI.5.6}
Note that to establish result 5.2, and consequently also the first assertion in 5.3, as well as 5.4, the reference to 4.2 can be replaced by a reference to VIII~6.4, whose proof is much easier. This also shows that the hypothesis $G$ affine over $S$ is unnecessary there. In addition, in \No6 we shall show how a variant of this method allows these results to be extended to the case of certain groups $H$ more general than groups of multiplicative type. These same observations extend to variants 5.2 bis etc. On the other hand, result 5.8 below uses essentially all the hypotheses made (in particular $G$ affine and smooth over $S$, $H$ of multiplicative type), and the lecturer knows no other proof of it than via the representability theorems 4.1 or 4.2{\renewcommand{\thefootnote}{*}\footnote{The situation has changed since the writing of this text, \cf XV and XIX \No6.}\addtocounter{footnote}{-1}}.
\end{remark}

\numpara{5.7}
Since the morphism (x), resp. (x bis), is smooth, hence open, its image is open. Let $R$ be this image, endowed with the structure induced by $M\times_S M$; one observes easily that $R$ is an \emph{equivalence relation} in $M$, which is moreover nothing other than that described explicitly in 5.4, resp. 5.4 bis. It would be interesting to know whether the quotient sheaf $M/R$ (which is formally étale over $S$) is representable (it is then representable by an étale prescheme over $S$); it is so, for example, if $S$ is the spectrum of a field. One will note, moreover, that it can happen that $R$ is not closed in $M\times_S M$ (even if $G$ is quasi-finite over $S\ldots$), which means (when $M/R$ is representable) that $M/R$ is then not separated over $S$.

\begin{theorem}{5.8}\label{XI.5.8}
Let $S$ be a prescheme, $G$ an $S$-prescheme smooth and affine over $S$, $s\in S$. Then:
% typo?: no group structure specified for G in the printed statement.
\begin{enumerate}[label=\alph*)]
\item For every subgroup $H_0$ of multiplicative type of $G_s$, there exists an étale morphism $S'\to S$, a point $s'$ of $S'$ above $s$ such that the residue extension $\kappa(s')/\kappa(s)$ is trivial, and a subgroup $H'$ of multiplicative type of $G'=G\times_S S'$ such that $H'_{s'}=H_0\otimes_{\kappa(s)}\kappa(s')$.
\item For every group homomorphism $u_0:H_s\to G_s$, where $H$ is an $S$-group prescheme of multiplicative type and of finite type, there exists an étale morphism $S'\to S$, a point $s'$ of $S'$ above $s$ such that the residue extension $\kappa(s')/\kappa(s)$ is trivial, and a group homomorphism $u':H'\to G'$ such that $u'_{s'}:H'_{s'}\to G'_{s'}$ is equal to $u_0\otimes_{\kappa(s)}\kappa(s')$.
\end{enumerate}
\end{theorem}
This follows from 4.1, resp. 4.2, and Hensel's lemma in the form 1.10.

\begin{proposition}{5.9}\label{XI.5.9}
Let $S$ be a prescheme, $G$ an $S$-prescheme smooth and affine over $S$, $H$ a subgroup of multiplicative type of $G$. Consider $\uCentr_G(H)\hookrightarrow\uNorm_G(H)$, which by 5.3 and 5.3 bis are closed group sub-preschemes of $G$, smooth over $S$. Then the first group is an open and closed sub-prescheme of the second, and the quotient sheaf
\[
W_G(H)=\uNorm_G(H)/\uCentr_G(H)
\]
is representable by an open group sub-prescheme of $\underline{\Aut}_{S\text{-gr}}(H)$; hence it is an $S$-group prescheme quasi-finite, étale and separated over $S$.
\end{proposition}
Indeed, consider the obvious homomorphism
\[
\theta:\uNorm_G(H)\to\underline{\Aut}_{S\text{-gr}}(H),
\]
whose kernel is by definition $\uCentr_G(H)$. Since $\underline{\Aut}_{S\text{-gr}}(H)$ is representable by an $S$-group prescheme étale and separated over $S$ (X~5.10), its identity section is an open and closed immersion, so its inverse image under $\theta$ is an open and closed subgroup of $\uNorm_G(H)$. I say, in addition, that $\theta$ is a \emph{smooth morphism}: this follows indeed formally from the definitions, and from the fact that $\uNorm_G(H)$ is smooth over $S$ and $\underline{\Aut}_{S\text{-gr}}(H)$ is étale over $S$. One concludes as in 5.3 that the image of $\theta$ is an open subset $U$ of $\underline{\Aut}_{S\text{-gr}}(H)$ and that, endowed with the induced structure, $U$ represents the quotient sheaf $\uNorm_G(H)/\uCentr_G(H)$. The latter is therefore étale and separated over $S$, since $\underline{\Aut}_{S\text{-gr}}(H)$ is, and it is quasi-finite over $S$, being quasi-compact over $S$ as the image of $\uNorm_G(H)$, which is so. This completes the proof of 5.9.
\begin{corollary}{5.10}\label{XI.5.10}
For every $s\in S$, let
\[
w(s)=\text{rank }\uNorm_{G_s}(H_s)/\uCentr_{G_s}(H_s)
\]
(which is also the index of $\Centr_{G(k)}H(k)$ in $\Norm_{G(k)}H(k)$, where $k$ is an algebraic closure of $\kappa(s)$). Then the function $s\mto w(s)$ is lower semicontinuous. For it to be constant near the point $s$, it is necessary and sufficient that $W_G(H)$ be finite over $S$ near $s$.
\end{corollary}
Indeed, for every $S$-prescheme $W$ that is étale, of finite type and separated over $S$, it is true that the function $s\mto w(s)=\text{rank }[W_s:\kappa(s)]$ is lower semicontinuous, and that it is constant near the point $s$ if and only if $W$ is finite over $S$ near $s$ (a fact pointed out in SGA~1, I~10.9, whose proof, which offers no difficulty, will be found in EGA~IV{\renewcommand{\thefootnote}{*}\footnote{EGA~IV~15.5.1 and 18.10.7.}\addtocounter{footnote}{-1}}).
% typo: closing parenthesis missing in the printed proof -> supplied.
\begin{remark}{5.11}\label{XI.5.11}
\nde{For a generalization to the non-affine case, see M.~Raynaud, \textit{Faisceaux amples sur les schémas en groupes et les espaces homogènes}, Lecture Notes Math.~119 (1970), IX.2.8.}Let $G$ be a group prescheme affine and smooth over $S$, $H$ a subgroup of multiplicative type; then 5.3 and 5.3 bis imply that the quotients
\[
G/\uCentr_G(H)\quad\text{and}\quad G/\uNorm_G(H)
\]
are smooth preschemes over $S$ and \emph{quasi-affine} over $S$, since in either case the moduli prescheme $M$ into which the quotient is embedded is such that every open subset $U$ of $M$ quasi-compact over $S$ is quasi-affine over $S$ (3.11). We shall moreover see in the following exposé\nde{\Cf remark XII.5.7. Let us mention here the following generalization. Without any affineness hypothesis on $G$, if $H$ is a torus and $S$ is locally noetherian, $U$ is affine and smooth by a theorem of Raynaud (\loccit, IX.2.6 and IX.2.8). Indeed, $G/\uNorm_G(H)$ is an open subset of $M$ and the largest representable open subset of $M$ is a disjoint sum of smooth affine open subsets over $S$.} that for every $U$ as above, the schematic closure $\overline U$ of $U$ in $M$ is even affine over $S$, which shows that the quotients considered are affine over $S$ provided the open subset $U$ in the corresponding moduli scheme (of group homomorphisms from $H$ to $G$, resp. of subgroups of multiplicative type of $G$) is closed in $M$. This is, for example, the case if $H$ is a ``maximal torus'' in the second case considered, or when $S$ is the spectrum of a field, \cf XII.5.4\nde{Under the hypothesis that $G$ has locally constant reductive rank.}. I do not know whether the quotients considered are affine over $S$ in general.
\end{remark}

\par\addvspace{1.1\baselineskip}
\noindent{\large\bfseries 6.\ On a rigidity property for homomorphisms of certain group schemes, and the representability of certain transporters{\renewcommand{\thefootnote}{*}\footnote{The present \No does not use the results of Nos.~3, 4, 5; its natural place would be in VI$_{\mathrm B}$.}\addtocounter{footnote}{-1}}}\par
\addvspace{0.45\baselineskip}
\addcontentsline{toc}{section}{6. On a rigidity property for homomorphisms of certain group schemes, and the representability of certain transporters}
\begin{proposition}{6.1}\label{XI.6.1}
Let $S$ be a locally noetherian prescheme, $H$ a commutative $S$-group prescheme of finite type over $S$, $E$ a set of integers $>0$ stable under multiplication, and suppose the following conditions hold:
\begin{enumerate}[label=\alph*)]
\item For every $n\in E$, the subgroup ${}_nH=\Ker(n\cdot\mathrm{id}_H)$ is finite and flat over $S$.
\item For every $s\in S$, the family of ${}_nH_s$ ($n\in E$) is schematically dense in $H$.
\end{enumerate}
Let $Y$ be a closed sub-prescheme of $H$. Then the subfunctor $\prod_{H/S}Y/H$ of $S$ (compare VIII~6) is representable by a closed sub-prescheme $T$ of $S$, and if $S$ is noetherian, there exists an $n\in E$ such that
\begin{equation}\tag{x}
\prod_{H/S}Y/H=\prod_{{}_nH/S}Y\cap{}_nH/S.
\end{equation}
% typo?: the denominator S on the right and H rather than H_s in b) kept as printed.
\end{proposition}
Indeed, by VIII~6.4, since by condition a) ${}_nH$ is finite and flat and a fortiori ``essentially free'' over $S$, it follows that the second member of (x) is representable by a closed sub-prescheme $T_n$ of $T$. Of course, for $n\in E$, $E$ ordered by divisibility, the $T_n$ form a decreasing family of closed sub-preschemes of $S$, so if $S$ is noetherian (which we can suppose), it is stationary for large $n$.
% typo?: T_n of T kept as printed.

Let $T$ be the value of $T_n$ for large $n$; I say that $T$ does indeed represent the first member of (x), which will complete the proof. One is reduced to proving that if $S'$ is a prescheme over $S$ such that $(Y\cap{}_nH)_{S'}=({}_nH)_{S'}$ for every $n\in E$, \ie such that $Y_{S'}\supset{}_nH_{S'}$ for every $n\in E$, then $Y_{S'}=H_{S'}$. Now this is indeed the case, for by IX~4.4 the family of ${}_nH_{S'}$ is schematically dense in $H_{S'}$, taking account of conditions a) and b).
\begin{theorem}{6.2}\label{XI.6.2}
Let $S$ be a prescheme, $H,E$ as in 6.1 satisfying conditions a), b), $u:H\to G$ a homomorphism of $S$-groups from $H$ to an $S$-group prescheme locally of finite type over $S$, and finally $K$ a closed group sub-prescheme of $H$. Consider the subfunctor $\uTransp_G(u,K)$ of $G$ (\cf 2.5). Then the latter is representable by a closed sub-prescheme of $G$, and if $G$ is noetherian (for example $S$ noetherian and $G$ of finite type over $S$), there exists an integer $n\in E$ such that $\uTransp_G(u,K)=\uTransp_G(u|_{{}_nH},K)$. If finally $G$ is smooth over $S$, and $K$ smooth over $S$ or of multiplicative type, and if $H,E$ satisfy the following condition, stronger than a):

 a') For every $n\in E$, the subgroup ${}_nH$ of $H$ is of multiplicative type,

then $\uTransp_G(u,K)$ is smooth over $S$.
% typo?: K a subgroup of H kept as printed.
\end{theorem}
Indeed, consider the $G$-group $H_G=H\times_S G$, which obviously satisfies the conditions of 6.1 ($S$ replaced by $G$, and $H$ by $H_G$), and the closed sub-prescheme $Y$ of $H_G$, inverse image of $K\subset G$ under
\[
H_G=H\times_S G\to G
\]
defined by $(h,g)\mto\operatorname{int}(g)\cdot u(h)$. Then $\uTransp_G(u,K)$ is nothing other than $\prod_{H_G/G}Y/H_G$ (compare VIII, examples 6.5 e)). Thus the first assertions follow from 6.1, and moreover one sees that for every quasi-compact open subset $U$ of $G$, there exists $n\in E$ such that $\uTransp_G(u,K)$ and $\uTransp_G(u|_{{}_nH},K)$ have the same trace on $U$. To verify the last assertion of 6.2, one can therefore replace $H$ by an ${}_nH$, and then it suffices to apply 2.5, which applies since ${}_nH$ is supposed of multiplicative type over $S$.
\begin{remark}{6.3}\label{XI.6.3}
The preceding proof uses only the very elementary result VIII~6.4, and in addition (for the last part) 2.5, that is, when $K$ is smooth over $S$, the infinitesimal result IX~3.6, hence the vanishing of the cohomology of groups of multiplicative type. One will note that in the most important cases (\cf 6.7) one can even suppose the $n\in E$ prime to the residue characteristics of $S$, \ie invertible in $\OS$, hence the ${}_nH$ finite étale over $S$, and then the cohomological result invoked is practically trivial, so 6.2 is then independent of the theory of groups of multiplicative type.
\end{remark}

\numpara{6.4}
One sees as usual that 6.1 extends to the case where one has an $S$-prescheme $S'$ over $S$ (not necessarily locally noetherian), and a closed sub-prescheme $Y'$ of $H'=H_{S'}$, provided $Y'\to H'$ is of finite presentation, \ie the ideal defining $Y'$ of finite type: then $\prod_{H'/S'}Y'/H'$ is representable by a closed sub-prescheme $T'$ of $S'$, such that $T'\to S'$ is of finite presentation, and if $S'$ is quasi-compact, there exists an $n\in E$ such that the relation analogous to (x) holds.
